Con motivo del Mundial de Fútbol 2026, la FIFA ha diseñado una red logística para transportar material deportivo, equipaciones, equipos de retransmisión y suministros médicos desde un centro logístico principal hasta el Centro de Operaciones de la Final.

La red se modela mediante un grafo dirigido en el que cada arco tiene asociada una capacidad máxima de transporte expresada en toneladas por día.

Los nodos de la red son:

\begin{itemize}
    \item $A$: Aeropuerto de Ciudad de México.
    \item $B$: Aeropuerto de Monterrey.
    \item $C$: Aeropuerto de Guadalajara.
    \item $D$: Centro de Distribución Norte.
    \item $E$: Centro de Distribución Centro.
    \item $F$: Centro de Distribución Sur.
    \item $G$: Estadio Azteca.
    \item $H$: Estadio BBVA.
    \item $I$: Estadio Akron.
    \item $S$: Centro Logístico Central.
    \item $T$: Centro de Operaciones de la Final .
\end{itemize}

Las capacidades de transporte de los arcos son las siguientes:

\[
\begin{array}{|c|c|}
\hline
\text{Arco} & \text{Capacidad} \\
\hline
S \rightarrow A & 30 \\
S \rightarrow B & 20 \\
S \rightarrow C & 25 \\
A \rightarrow D & 15 \\
A \rightarrow E & 10 \\
B \rightarrow D & 10 \\
B \rightarrow F & 15 \\
C \rightarrow E & 15 \\
C \rightarrow F & 10 \\
D \rightarrow G & 20 \\
D \rightarrow H & 5 \\
E \rightarrow G & 10 \\
E \rightarrow I & 15 \\
F \rightarrow H & 15 \\
F \rightarrow I & 10 \\
G \rightarrow T & 20 \\
H \rightarrow T & 15 \\
I \rightarrow T & 20 \\
\hline
\end{array}
\]

\begin{enumerate}
    \item Representa la red mediante un grafo dirigido indicando las capacidades en cada arco.

    \item Determina el flujo máximo que puede enviarse desde el nodo  $S$ hasta el nodo  $T$ utilizando el algoritmo de Ford-Fulkerson; obteniendo una distribución de flujo máxima para todos los arcos de la red.

    \item Determina un corte mínimo de la red y calcula su capacidad.

    \item Interpreta el significado del corte mínimo en el contexto de la logística del Mundial de Fútbol 2026.
\end{enumerate}